\begin{lemma}[Scheide maximal-tree package (b)]
Let $(G,e_0,\varphi)$ be a $k$-triple and let $T$ be a fixed-pair maximal
Tashkinov tree for $(e_0,\varphi)$.  Then $|V(T)|$ is odd.
\end{lemma}
\begin{proof}
Put $X=V(T)$.  The positive-length condition in \noderef{TreeSequence}
makes $X$ nonempty.  By \noderef{ScheideMaximalElementaryClosed}, $X$ is
elementary and closed.  Every vertex of $X$ misses a color.  Indeed, the
$k$-triple inequality is $k\geq\Delta(G)+1$ by \noderef{KTriple}; the
number of incident edge identities is the degree from \noderef{Degree},
and \noderef{MaximumDegree} bounds this degree by $\Delta(G)$.  Properness
makes their present colors distinct, so fewer than $k$
colors are present at every vertex.  The complement definition
\noderef{MissingColors} therefore supplies a vertex $x\in X$ and a color
$\eta$ missing at $x$.

Elementariness says that $\eta$ is not missing, hence is present, at every
vertex of $X\setminus\{x\}$.  By \noderef{PresentColors}, each such vertex
is incident with an $\eta$-colored edge identity.  Properness gives at most
one such edge at each vertex.  Closedness says that no $\eta$-colored edge
has exactly one end in $X$, because $\eta$ is missing on $X$.  Consequently
the $\eta$-colored internal edges form a matching which saturates exactly
the $|X|-1$ vertices other than $x$.  Each matching edge has two distinct
ends by \noderef{LooplessMultigraph}, so $|X|-1$ is twice the number of
these edge identities and is even.  Hence $|X|$ is odd.
\end{proof}
